\chapter{Core Definitions}
\label{ch:definitions}

\section{Motivation}

We want a framework in which \emph{addition} and \emph{scaling} make sense and
nothing else is assumed. The gain is uniformity: one theorem proved in that
framework applies simultaneously to tuples of numbers, to polynomials, to
matrices and to functions.

\section{Vector spaces}

\begin{definition}{Vector space}{vector-space}
Let $\K$ be a field. A \term{vector space} over $\K$ is a set $E$ equipped with
an addition $E\times E\to E$ and a scalar multiplication $\K\times E\to E$ such
that $(E,+)$ is an abelian group and, for all $\lambda,\mu\in\K$ and
$x,y\in E$,
\[
  \lambda(x+y)=\lambda x+\lambda y,\qquad
  (\lambda+\mu)x=\lambda x+\mu x,\qquad
  (\lambda\mu)x=\lambda(\mu x),\qquad
  1\cdot x=x.
\]
The elements of $E$ are called vectors, those of $\K$ scalars.
\end{definition}

\begin{example}{The four standard examples}{standard-spaces}
\begin{enumerate}
  \item $\K^{n}$ with componentwise operations.
  \item $\K[X]$, and its subspace $\K_{n}[X]$ of polynomials of degree at most $n$.
  \item $\Mat{m,n}{\K}$, the $m\times n$ matrices.
  \item $\mathcal{F}(X,\K)$, all maps from a set $X$ to $\K$.
\end{enumerate}
Every example met in this book is one of these four, or a subspace of one of them.
\end{example}

\begin{example}{Three non-examples}{non-examples}
\begin{enumerate}
  \item $\N^{2}$ is not a vector space over $\R$: it is not stable under
        multiplication by $-1$.
  \item $\set{(x,y)\in\R^{2}\setst xy=0}$ is stable under scaling but not under
        addition: $(1,0)+(0,1)=(1,1)$.
  \item The set of polynomials of degree \emph{exactly} $n$ is not a subspace:
        it does not contain $0$, and $X^{n}+(-X^{n}+1)$ has degree $0$.
\end{enumerate}
\end{example}

\begin{definition}{Subspace}{subspace}
A subset $F\subseteq E$ is a \term{linear subspace} if $0\in F$ and
$\lambda x+\mu y\in F$ for all $x,y\in F$ and $\lambda,\mu\in\K$. A subspace is
itself a vector space for the restricted operations.
\end{definition}

\begin{proposition}{Intersections and sums}{sum-intersection}
If $F$ and $G$ are subspaces of $E$, then $F\cap G$ and
$F+G=\set{x+y\setst x\in F,\ y\in G}$ are subspaces, and $F+G$ is the smallest
subspace containing $F\cup G$. In general $F\cup G$ is \emph{not} a subspace.
\end{proposition}

\begin{proof}
Stability of $F\cap G$ and of $F+G$ is immediate. Any subspace $H$ containing
$F\cup G$ contains every sum $x+y$ with $x\in F$, $y\in G$, hence contains
$F+G$; and $F+G$ contains $F\cup G$. The last assertion is \cref{exo:union}.
\end{proof}

\begin{definition}{Direct sum}{direct-sum}
The sum $F+G$ is \term{direct}, written $F\oplus G$, if $F\cap G=\set{0}$.
Equivalently, every $x\in F+G$ has a unique decomposition $x=y+z$ with $y\in F$
and $z\in G$.
\end{definition}

\begin{intuition}
A direct sum is a coordinate system: to say $E=F\oplus G$ is to say that each
vector of $E$ has exactly one address, one component in $F$ and one in $G$.
Uniqueness is what makes projections well defined.
\end{intuition}

\section{Families of vectors}

\begin{definition}{Linear combination, span}{span}
Let $(x_i)_{i\in I}$ be a family of vectors of $E$. Its \term{span}, written
$\Span(x_i)_{i\in I}$, is the set of finite linear combinations
$\sum_{i\in J}\lambda_i x_i$ with $J\subseteq I$ finite and $\lambda_i\in\K$. It
is the smallest subspace containing the family.
\end{definition}

\begin{definition}{Free family, generating family, basis}{basis}
A family $(x_i)_{i\in I}$ is
\begin{itemize}
  \item \term{free} (or linearly independent) if every finite relation
        $\sum_{i\in J}\lambda_i x_i=0$ forces $\lambda_i=0$ for all $i\in J$;
  \item \term{generating} if $\Span(x_i)_{i\in I}=E$;
  \item a \term{basis} if it is both free and generating.
\end{itemize}
\end{definition}

\begin{proposition}{Characterisation of a basis}{basis-unique}
A family $\mathcal{B}=(e_i)_{i\in I}$ is a basis of $E$ if and only if every
vector of $E$ is a linear combination of $\mathcal{B}$ in exactly one way.
\end{proposition}

\begin{proof}
If $\mathcal{B}$ is a basis, existence follows from generation, and if
$\sum\lambda_i e_i=\sum\mu_i e_i$ then $\sum(\lambda_i-\mu_i)e_i=0$, so
$\lambda_i=\mu_i$ by freeness. Conversely, existence of a decomposition gives
generation, and uniqueness applied to $0=\sum 0\cdot e_i$ gives freeness.
\end{proof}

\begin{definition}{Dimension}{dimension}
If $E$ admits a finite basis, all its bases have the same cardinality
(\cref{thm:dimension-well-defined}); this common number is the \term{dimension}
$\dim E$. Otherwise $E$ is said to be infinite-dimensional.
\end{definition}

\section{Linear maps}

\begin{definition}{Linear map}{linear-map}
A map $u\colon E\To F$ is \term{linear} if $u(\lambda x+\mu y)=\lambda u(x)+\mu u(y)$
for all $x,y\in E$ and $\lambda,\mu\in\K$. The set of such maps is written
$\Hom(E,F)$, and $\End(E)=\Hom(E,E)$. An \term{endomorphism} is an element of
$\End(E)$, an \term{isomorphism} is a bijective linear map.
\end{definition}

\begin{definition}{Kernel, image, rank}{kernel-image}
For $u\in\Hom(E,F)$ one sets
\[
  \Ker u=\set{x\in E\setst u(x)=0},\qquad
  \im u=\set{u(x)\setst x\in E},\qquad
  \rank u=\dim\im u .
\]
$\Ker u$ is a subspace of $E$ and $\im u$ a subspace of $F$.
\end{definition}

\begin{proposition}{Injectivity and the kernel}{injective-kernel}
A linear map $u$ is injective if and only if $\Ker u=\set{0}$.
\end{proposition}

\begin{proof}
If $u$ is injective and $u(x)=0=u(0)$ then $x=0$. Conversely, if
$\Ker u=\set{0}$ and $u(x)=u(y)$, then $u(x-y)=0$, so $x-y=0$.
\end{proof}

\begin{definition}{Projection}{projection}
An endomorphism $p\in\End(E)$ with $p^{2}=p$ is called a \term{projection}.
\end{definition}

\Cref{fig:projection} shows the geometric content of the identity $p^2=p$:
projecting a second time changes nothing.

\begin{figure}[htbp]
  \centering
  % Orthogonal projection onto a line: the geometric picture behind p^2 = p.
\begin{tikzpicture}[scale=1.05,>=Stealth,font=\small]
  \draw[dmzrule,line width=0.4pt] (-0.6,0) -- (4.4,0);
  \draw[dmzrule,line width=0.4pt] (0,-0.6) -- (0,3.1);
  % the line D = im p
  \draw[dmzaccent,line width=1pt] (-0.4,-0.2) -- (4.2,2.1) node[right,text=dmzaccent] {$D=\im p$};
  % vector x and its projection
  \draw[->,dmzink,line width=0.9pt] (0,0) -- (1.6,2.6) node[above left] {$x$};
  \draw[->,dmztheorem,line width=0.9pt] (0,0) -- (2.65,1.33) node[below right,text=dmztheorem] {$p(x)$};
  \draw[dmzsoft,dashed] (1.6,2.6) -- (2.65,1.33);
  \draw[dmzsoft,line width=0.4pt] (2.65,1.33) ++(-0.18,0.09) -- ++(0.09,0.18) -- ++(0.18,-0.09);
  % the kernel direction
  \draw[->,dmzgreen,line width=0.8pt] (0,0) -- (-1.05,1.27) node[above,text=dmzgreen] {$\Ker p$};
  \node[dmzsoft,anchor=north west,font=\footnotesize,text=dmzsoft] at (0.15,-0.15) {$0$};
\end{tikzpicture}

  \caption{A projection $p$ decomposes the plane as $\Ker p\oplus\im p$.}
  \label{fig:projection}
\end{figure}

\section{Eigenvalues}

\begin{definition}{Eigenvalue, eigenvector, spectrum}{eigenvalue}
Let $u\in\End(E)$. A scalar $\lambda\in\K$ is an \term{eigenvalue} of $u$ if
$\Ker(u-\lambda\id)\neq\set{0}$; a nonzero element of that kernel is an
\term{eigenvector}, and $E_\lambda=\Ker(u-\lambda\id)$ is the associated
\term{eigenspace}. The set of eigenvalues is the \term{spectrum} $\spec(u)$.
\end{definition}

\begin{definition}{Diagonalisable endomorphism}{diagonalisable}
An endomorphism $u$ of a finite-dimensional space $E$ is \term{diagonalisable}
if $E$ admits a basis of eigenvectors of $u$; equivalently, if
$E=\bigoplus_{\lambda\in\spec(u)}E_\lambda$.
\end{definition}

\begin{definition}{Annihilating polynomial}{annihilating}
A polynomial $P\in\K[X]$ \term{annihilates} $u$ if $P(u)=0$. The monic
generator of the ideal of annihilating polynomials is the \term{minimal
polynomial} $\pi_u$.
\end{definition}

\begin{warning}
A definition is not a criterion. \Cref{def:diagonalisable} says what
diagonalisable means; it gives no method. The usable criteria are
\cref{thm:diagonalisable-criterion} and its corollary on split annihilating
polynomials with simple roots.
\end{warning}
